\section{The \instname\ data reduction flow.}
\label{sec:reduction_flow}


The overall data flow of the \instname\ pipeline is displayed in Figure \ref{fig:reduction_cascade}.

 
\begin{figure}[h!]
\centering
\includegraphics[width=0.9\textwidth]{../instrument_tutorial/figures/reduction_cascade.jpg}
\caption{The data reduction cascade of the \instname\ workflow.}
\label{fig:reduction_cascade}
\end{figure}



The reduction cascade is organized in tasks, which represent
well-defined steps in the process. Tasks can be grouped inside
sub-workflows.  Each task runs a recipe; the detailed description of
the algorithms, input, outputs and recipe parameters used in each
recipe are available in the pipeline manual. Here, we present only the
description of most important features.

The \wkfname\  \edps\ workflow is designed to execute the tasks that
deliver the final reduced data cube for each dataset. It can be either
the product of a single exposure, or the combination of multiple
exposures. Only calibrations needed by the selected the scientific
exposures are processed.

It is possible to set \edps\ to perform the data reduction until a
certain step of the reduction chain (e.g. to reduce only standar
stars, or only flat fields).  This is done by specifying the desired
tasks in the field {\bf Select reduction target} of the {\bf Raw Data}
tab.

The reduction steps of the \wkfname\ workflow are listed below. Before
starting the reduction, the parameters of the recipes associated to
each task can be configured by pressing the button

\includegraphics[width=0.8cm]{../common/figures/configure_dataset.jpg}
close to each dataset configuration.
See for more info on the configuration editor \ref{sec:configuration_editor}



UVES can be operated standalone or together with the FLAMES instrument. If a
destinction between these two operational modes is needed in this tutorial, then
the former is referred to as `UVES (Slit)' and the latter as `UVES-Fibre'. 

The two spectroscopic arms of UVES, i.e., the blue and the red arm, are processed separately. 
For the blue arm, the pipeline recipes create one product per product category whereas for the
red arm, usually two products are created, one for the `lower' and one for the `upper' chip. These
different products can be distinguished by the tags {\tt \_BLUE}, {\tt \_REDL}, and {\tt \_REDU} in their
{\tt PRO CATG} header keywords. For example, a complete set of raw blue and red bias calibrations
gives three master calibrations having {\tt PRO CATG} values of {\tt MASTER\_BIAS\_BLUE}, {\tt MASTER\_BIAS\_REDL}, and
{\tt MASTER\_BIAS\_REDU}, respectively. In the following description, {\tt <CHIP>} can be one of
{\tt BLUE}, {\tt REDL}, or {\tt REDU}.

\section{UVES (Slit) data reduction}
\subsection{Master bias}

Task: {\tt bias}. Recipe: {\tt uves\_cal\_mbias}. Output: {\tt MASTER\_BIAS\_<CHIP>}.

This task creates a master bias for each chip. A minimum of three input frames is required.

\subsection{Initial order and wavelength solution}

Task: {\tt predict}. Recipe: {\tt uves\_cal\_predict}. Output: {\tt LINE\_GUESS\_TAB\_<CHIP>} and \newline
{\tt ORDER\_GUESS\_TAB\_<CHIP>}.

This task creates the initial guess for the order and wavelength identification.

\subsection{Order positions}

Task: {\tt orderpos}. Recipe: {\tt uves\_cal\_orderpos}. Output: {\tt ORDER\_TABLE\_<CHIP>}.

This task creates the final identification of the Echelle orders.

\subsection{Master flat field}

Task: {\tt flat}. Recipe: {\tt uves\_cal\_mflat}. Main output: {\tt  MASTER\_FLAT\_<CHIP>}.

This task creates the master flat for the `regular' lamp flat fields (which have {\tt DPR TYPE = 
LAMP,FLAT}). The workflow also supports the usage of flat fields taken with the deuterium lamp
({\tt DPR TYPE = LAMP,DFLAT}) and of flats where the iodine cell is inserted into the light path
({\tt DPR TYPE = LAMP,IFLAT}). 

The deuterium lamp flats are recommended for the spectral region shorter than 350 nm. 
They are taken as supplementary calibrations for
the UVES blue arm central wavelength setting of 346 nm. 
If deuterium flats are present in the data set,
then they are combined with the regular lamp flats into a combined master flat that provides much better
S/N shorter than 350 nm. The tasks {\tt dflat} and {\tt combined\_flat} are executed in addition to the
{\tt flat} task.

The iodine cell flats are provided only for the red 600 nm setting. They contain the the iodine cell
absorption lines superimposed on the flat-field spectrum. The task {\tt iflat} is executed then.

\subsection{Wavelength calibration}

Task: {\tt arc}. Recipe: {\tt uves\_cal\_wavecal}. Output: {\tt LINE\_TABLE\_<CHIP>}.

The wavelength calibration is determined from Th-Ar arc-lamp exposures.

\subsection{Response curve}

Task: {\tt response}. Recipe: {\tt uves\_cal\_response}. Output: {\tt INSTR\_RESPONSE\_FINE\_<CHIP>}.

Flux calibration can be obtained either by processing an observation of a flux standard or by
applying a master response curve. If a raw standard star calibration is present in the data set,
then this observation is processed by the {\tt response} task. The resulting response curve is then
used during processing of the science data. Master response curves are available for a number of
setups. If a raw standard star observation is not present but a matching master response is found, then
the latter is used for the flux calibration. If neither a raw standard nor a master response is 
present in the data set, then the extracted science spectrum is not flux calibrated.

\subsection{Spectrum extraction}
\label{sec:extraction}

Task: {\tt object}. Recipe: {\tt uves\_obs\_scired}. 

This task extracts the spectrum and wavelength calibrates it. If a response curve is available, then 
the spectrum is also flux-calibrated. The main output files are listed in Table~\ref{tab:object-output}.

\begin{table}[htbp]
\centering
\begin{tabular}{|l|l|}
\hline
\textbf{Product} & \textbf{Description} \\
\hline 
{\tt RED\_SCI\_<TYPE>\_<CHIP>}               &extracted, flat-fielded, wavelength-calibrated, \\
                                             &order-merged, and sky-subtracted science spectrum\\
{\tt ERRORBAR\_SCI\_<TYPE>\_<CHIP>}          &uncertainty on {\tt RED\_SCI\_<TYPE>\_<CHIP>}\\
{\tt FLUXCAL\_SCI\_<TYPE>\_<CHIP>}           &flux-calibrated version of\\
                                             &{\tt RED\_SCI\_<TYPE>\_<CHIP>}\\
{\tt FLUXCAL\_ERRORBAR\_SCI\_<TYPE>\_<CHIP>} &uncertainty on {\tt FLUXCAL\_SCI\_<TYPE>\_<CHIP>}\\
\hline
\end{tabular}
\caption{Main output products from the {\tt object} task. For the values of {\tt <TYPE>}, see Table~\ref{tab:object-types}}
\label{tab:object-output}
\end{table}

The {\tt PRO CATG} depends on whether or not the image slicer and/or the absorption cell has been used.
See Table~\ref{tab:object-types}. 

\begin{table}[htbp]
\centering
\begin{tabular}{|l|l|}
\hline
\textbf{<TYPE>} & \textbf{Description} \\
\hline 
{\tt POINT}               &Regular slit observation\\
{\tt SLICER}              &Observation using the image slicer\\
{\tt ABSCELL}             &Observation using the absorption cell\\
{\tt SLICABSCELL}         &Observation using the image slicer and the absorption cell\\
\hline
\end{tabular}
\caption{Possible values for {\tt <TYPE>}}
\label{tab:object-types}
\end{table}

The object spectrum is extracted using either {\tt optimal} extraction (default), 
{\tt linear} extraction (summing the flux within the extraction slit), 
{\tt average} extraction (averaging the flux within the extraction slit),
or the {\tt 2d} method (resampling the full slit length to a 2-dimensional image,
with wavelength along the x-axis and position along the slit along the y-axis, 
best suited for extended sources). In the latter case, output products have an additional tag
`{\tt 2D}' in their {\tt PRO CATG} values. The recipe parameter that controls the extraction
method is {\tt reduce.extract.method}.
More details on spectrum extraction can be found in Sect.~\ref{sect:optimizing_extraction}.

If the spectrum is optimally extracted, cosmic rays are detected and rejected, 
and the sky is fitted and subtracted during the extraction. 
If linear/average extraction is used, then the sky is extracted in windows above and below
the object and subsequently subtracted 
(if the recipe parameter {\tt reduce.skysub} is set to the default value TRUE).

\subsection{Phase 3 compatible science product}

Task: {\tt science\_slit}. Recipe: {\tt uves\_utl\_idp}. \newline
Output: {\tt RED\_SCI\_<TYPE>\_BLUE} or {\tt RED\_SCI\_<TYPE>\_RED} or 
{\tt FLUXCAL\_SCI\_<TYPE>\_BLUE} or {\tt FLUXCAL\_SCI\_<TYPE>\_RED}.

This task takes the output from {\tt object} and converts it into a single product that is
compatible with the ESO Phase 3 standard. The files are converted from
1D image format into a table with columns for the wavelength, the flux, the error, etc. Also,
the two spectra for the red arm are merged into a single file.

\subsection{Spectrum combination}

Task: {\tt combine\_spectra}. Recipe: {\tt esotk\_spectrum1d\_combine}. \newline
Output: {\tt ESOTK\_SPECTRUM\_IDP\_FORMAT}.

This task stacks two or more spectra from the same target (and with the same setup).
It requires as input the output from {\tt science\_slit}.

\section{UVES-Fibre data reduction}

In case of UVES-Fibre, <CHIP> can be `REDL' or `REDU'.

\subsection{Master bias}

Task: {\tt bias}. Recipe: {\tt uves\_cal\_mbias}. Output: {\tt MASTER\_BIAS\_<CHIP>}.

This task creates a master bias for each chip. A minimum of three input frames is required.

\subsection{Initial order and wavelength solution}

Task: {\tt predict\_mos}. Recipe: {\tt flames\_cal\_predict}.
\newline Output: {\tt FIB\_LIN\_GUE\_<CHIP>} and 
{\tt FIB\_ORD\_GUE\_<CHIP>}.

This task creates the initial guess for the order and wavelength identification.

\subsection{Order positions}

Task: {\tt orderpos\_mos}. Recipe: {\tt flames\_cal\_orderpos}. 
\newline Output: {\tt FIB\_ORD\_TAB\_<CHIP>} and {\tt FIB\_ORDEF\_<CHIP>}.

This task creates the final identification of the Echelle orders.

\subsection{Master flat field}

Task: {\tt flat\_mos}. Recipe: {\tt flames\_cal\_mkmaster}. Main output: {\tt  MASTER\_SFLAT\_<CHIP>}.

This task creates the (slit) master flat fields for three different positions on
the chip. Therefore, there are three files for each category.

\subsection{Fibre flat field}

Task: {\tt sff\_ofpos\_mos}. Recipe: {\tt flames\_cal\_prep\_sff\_ofpos}.\newline
Main output: {\tt FIB\_ORDEF\_TABLE\_<CHIP>}.

This task is used to determine the fibre order table and to construct several 
master calibrations needed to extract the science fibres.

\subsection{Wavelength calibration}

Task: {\tt arc\_mos}. Recipe: {\tt flames\_cal\_wavecal}. Output: {\tt FIB\_LINE\_TABLE\_<CHIP>}.

The wavelength calibration is determined from Th-Ar arc-lamp exposures.

\subsection{Spectrum extraction}

Task: {\tt object\_mos}. Recipe: {\tt flames\_obs\_scired}. 

This task extracts the spectra for each fibre. The main output files are listed
in Table~\ref{tab:mos_output}.

\begin{table}[htbp]
\centering
\begin{tabular}{|l|l|}
\hline
\textbf{Product} & \textbf{Description} \\
\hline 
{\tt MWXB\_SCI\_<CHIP>}               &extracted, flat-fielded, wavelength-calibrated, order-merged,\\
                                             &and fibre throughput corrected science spectrum;\\
                                      &one file for each fibre\\
{\tt ERR\_MWXB\_SCI\_<CHIP>}          &error for {\tt MWXB\_SCI\_<CHIP>}\\
\hline
\end{tabular}
\caption{Main output products from the {\tt object\_mos} task}
\label{tab:mos_output}
\end{table}